\section{Khoảng trống nghiên cứu}

Mặc dù kiến trúc truy xuất tăng cường sinh nội dung cơ bản đã chứng minh tính hiệu quả vượt trội trong việc cải thiện độ chính xác của các mô hình ngôn ngữ lớn, việc ứng dụng RAG vào tập dữ liệu về các bệnh lý về xương vẫn đối mặt với những rào cản kỹ thuật và thách thức nghiệp vụ lớn cần được giải quyết.

Khoảng trống thứ nhất xuất phát từ rào cản đứt gãy ngữ cảnh trong quá trình truy xuất tri thức y khoa. Tài liệu chuyên khoa về bệnh lý về xương có đặc thù chứa đựng mật độ thông tin cao và sự liên kết chéo phức tạp giữa các khái niệm y học (như mối quan hệ giữa triệu chứng lâm sàng, phác đồ phân bậc điều trị, và các chỉ định hay chống chỉ định đối với bệnh nhân có bệnh nền). Khi áp dụng các kỹ thuật RAG truyền thống dựa trên việc chia nhỏ văn bản và truy xuất biểu diễn vector đơn thuần, hệ thống dễ trích xuất các đoạn văn bản rời rạc, làm mất đi tính toàn vẹn của ngữ cảnh gốc. Sự đứt gãy này khiến mô hình ngôn ngữ lớn tiếp nhận ngữ cảnh thiếu sót, dẫn đến nguy cơ sinh ra câu trả lời không đầy đủ hoặc sai lệch. Các hướng tiếp cận hiện đại trong tổng quan tài liệu như truy xuất dựa trên đồ thị tri thức (GraphRAG) \cite{microsoft2024} hoặc các kỹ thuật truy xuất kết hợp từ khóa và ngữ nghĩa đang được chú ý khảo sát nhằm khắc phục giới hạn này.

Khoảng trống thứ hai nằm ở yêu cầu kiểm định độ tin cậy nghiêm ngặt và cơ chế minh bạch bằng chứng trong tư vấn y tế. Khác với các hệ thống hỏi đáp chung, một ứng dụng y khoa đòi hỏi từng luận điểm trả lời phải được chứng thực trực tiếp và trỏ đúng về văn bản tham chiếu chuẩn xác \cite{medtrustrag, imedrag}. Việc duy trì độ trung thực tuyệt đối so với nguồn tài liệu gốc, tránh sự suy diễn tự do của mô hình trong quá trình tổng hợp thông tin về các bệnh lý về xương là một bài toán kỹ thuật phức tạp, đòi hỏi phải có sự khảo sát thực nghiệm và đánh giá hệ thống trên các tiêu chí chuẩn hóa.
